% !TeX root = ../main.tex
\section{Products and open-source tools}
\label{sec:products}

Statements in this section come from manuals and product pages; vendors publish
no accuracy figures for recording-based tone features. Products take three kinds
of input (Table~\ref{tab:products}).

\paragraph{Capture of physical equipment.} Kemper profiling plays test signals
into the amplifier and analyses its response; compressors and time-based effects
must be bypassed~\cite{kemper2018manual}. Neural DSP's Neural Capture, IK
Multimedia's TONEX (with parametric models announced in October 2026), Line 6's
Proxy and NAM work the same way with neural
models~\cite{neuraldsp2026qcmanual,ik2026parametric,line6ndproxy,atkinsonndnam}.
None accepts a recording. The Quad Cortex manual says that a capture's gain
metadata is estimated from the device's dynamic range and harmonic content, the
only vendor statement found of a gain descriptor computed from
audio~\cite{neuraldsp2026qcmanual}.

\paragraph{Spectral matching to a recording.} The Axe-Fx III \textbf{Tone Match}
block computes long-term average spectra of a reference and of the user's preset
and applies their smoothed difference as a filter, exportable as a cabinet
IR~\cite{fractalndtonematch}. Its manual lists the conditions: a starting preset
with the right amplifier and the right amount of distortion, no effects, an
isolated and unclipped guitar track, matched loudness, mono analysis, broad
playing material. iZotope's Match EQ does the same for mixes~\cite{izotope2019matcheq}.
These tools estimate only the linear part, after the user has fixed the
nonlinear one. The app's final step is the same operation
(Section~\ref{sec:impl-match}), applied after its own search has fixed the gain.

\paragraph{Presets from audio, song names or text.} Positive Grid's BIAS X
``Music-to-Tone'' accepts a file of up to 15 minutes, full mixes included, finds a
guitar segment and returns a full chain; it runs online and no method or accuracy
is published~\cite{positivegridndbiasxhelp,positivegrid2025biasxai}. Small apps
generate presets from a song name with a language model's knowledge of the
artist's gear, sometimes with stem separation and coarse spectral
features~\cite{bouharb2026gptonebuilder,aitonetechndgenerator}; none publishes an
evaluation. The most detailed public design is a Chinese patent granted in 2025:
separation, a tone-type classifier, a network that picks an amp and cabinet model
and regresses five knobs, a fit of up to 24 peaking filters to the difference of
the spectra, tempo and key analysis for delay and modulation times, and retrieval
from a preset database for the remaining effects~\cite{huanyin2025presetpatent}.

\paragraph{Open-source tools.} NAM-EQ-Match renders the player's DI through every
capture in a NAM library at several input gains, compares EQ-invariant dynamics
features (crest factor, compression, spectral flux, MFCC statistics) with the
target, then fits a matching IR or tone-stack EQ~\cite{sd4562026nameqmatch}.
\code{tonematch} predicts seven amp knobs with a mel-spectrogram CNN and refines
them by CMA-ES through a grey-box amp, reporting a median log-spectral error of
1.54\dB{} on 24 isolated Guitar-TECHS cases against 3.68\dB{} with every knob at
noon~\cite{neuwirth2026tonematch}. ToneRunner retrieves plugin chains with a
contrastive encoder and documents how the encoder learns the riff instead of the
tone when the same dry clip is not shared across chains in
training~\cite{nathanben2026tonerunner}. These projects are small and not peer
reviewed, but they converge on the same design as the research: retrieve or
predict a nonlinear core, refine by rendering, finish with a linear match.

\begin{table}[h]
\centering
\caption{What products and tools take as input and return. ``Full chain'':
amplifier, cabinet and effects with settings. $^*$Method undocumented.}
\label{tab:products}
\small
\begin{tabular}{@{}lll@{}}
\toprule
System & Input & Output \\
\midrule
Capture products, NAM & the device itself & capture \\
Tone Match~\cite{fractalndtonematch} & isolated track & EQ / IR \\
BIAS X~\cite{positivegridndbiasxhelp} & recording, even a mix & full chain$^*$ \\
Song-name apps~\cite{bouharb2026gptonebuilder} & name, text & preset \\
CN118865922B~\cite{huanyin2025presetpatent} & recording & full chain \\
NAM-EQ-Match~\cite{sd4562026nameqmatch} & stem + player's DI & capture, EQ/IR \\
\code{tonematch}~\cite{neuwirth2026tonematch} & recording + DI & 7 amp knobs \\
Guitar Virtual Amp (this work) & separated stem + notes & app preset + match IR \\
\bottomrule
\end{tabular}
\end{table}
